% !TeX root = ../抽象代数笔记.tex
% 由主文档通过 \input 载入；请编译主文档。
\setcounter{chapter}{16} % 保留原书章号。
\chapter{有限域}\label{ch:finite-fields}

\begin{description}
  \item[学习目标] 掌握有限域的存在、唯一性、子域和自同构，能够实际构造并计算有限域。
  \item[先修知识] 第\ref{ch:splitting}章与循环群基础。
\end{description}

\section{有限域的阶与构造}
若有限域特征为素数 $p$，将其视作 $\mathbb F_p$ 上的有限维向量空间，证明其阶为 $p^n$，其中 $n\geq1$。反过来，在 $x^{p^n}-x$ 的分裂域中证明全部根组成一个含 $p^n$ 个元素的域，并证明此阶的域在同构意义下唯一，记为 $\mathbb F_{p^n}$。

\section{有限域乘法群与本原元素}
证明任意域的有限乘法子群为循环群，推出 $\mathbb F_q^\times$ 为 $q-1$ 阶循环群，其中 $q$ 为素数幂。定义有限域的本原元素（primitive element）为其乘法群生成元，并与第\ref{ch:splitting}章“生成整个扩张的本原元”区分；有限域中前者能生成域扩张，后者未必生成乘法群。

\section{弗罗贝尼乌斯自同构与子域}
在 $\mathbb F_{p^n}$ 上定义弗罗贝尼乌斯自同构（Frobenius automorphism）$\sigma(a)=a^p$，证明其阶为 $n$。证明 $\mathbb F_{p^n}$ 含有 $p^d$ 阶子域当且仅当 $d\mid n$，且对每个这样的 $d$ 子域唯一，等于 $x^{p^d}-x$ 的根集。以此显式计算全部固定 $\mathbb F_p$ 的自同构。

\section{不可约多项式与有限域运算}
对首一不可约多项式 $f\in\mathbb F_p[x]$、$\deg f=n$，把 $\mathbb F_p[x]/(f)$ 的元素表示成次数小于 $n$ 的剩余类代表元。给出加法、乘法与利用扩展欧几里得算法求逆的步骤。说明首一不可约多项式 $g\in\mathbb F_p[x]$ 整除 $x^{p^n}-x$ 当且仅当 $\deg g\mid n$。

\section{例题、练习与自检}
用不可约三次多项式构造 $\mathbb F_8$，列出其非零元素的幂，计算乘法逆元。列出 $\mathbb F_{p^{12}}$ 的全部子域并判断包含关系。练习中特别区分 $\mathbb F_{p^n}$ 与 $\mathbb Z/p^n\mathbb Z$：当 $n>1$ 时后者不是域。

